\documentclass[12pt]{article}
\usepackage{geometry}
\geometry{margin=1in}

% ----------------------------------------------------------
%  Linux Penguin Theme — Based on Recommended Colors
% ----------------------------------------------------------

% --- COLOR SETUP ---
\usepackage{xcolor}

\definecolor{cream}{RGB}{246,242,219}
\definecolor{teamred}{RGB}{209,57,64}
\definecolor{navy}{RGB}{40,46,87}
\definecolor{softwhite}{RGB}{250,250,230}
\definecolor{accentyellow}{RGB}{240,200,60}

% Extra utility colors
\definecolor{muted}{RGB}{120,125,140}
\definecolor{darkcream}{RGB}{230,225,205}

% ----------------------------------------------------------
%  FONT SETUP (Optional but Recommended)
%  Works best with XeLaTeX or LuaLaTeX
% ----------------------------------------------------------
\usepackage{fontspec}

% Main document font (Google fonts — upload to Overleaf)
% If you don't upload fonts, comment these two lines.
\setmainfont{SourceSans3}[
  Path=./,
  ItalicFont  = *-Italic,
  BoldFont    = *-Bold
]

\newfontfamily\headingfont{ScienceGothic}[
  Path=./,
]

% ----------------------------------------------------------
%  SECTION STYLE
% ----------------------------------------------------------
\usepackage{titlesec}

% Section
\titleformat{\section}
  {\Large\headingfont\bfseries\color{navy}}
  {\thesection}
  {1em}
  {\rule{\linewidth}{1pt}\vspace{0.5em}\\}

% Subsection
\titleformat{\subsection}
  {\large\headingfont\bfseries\color{teamred}}
  {\thesubsection}
  {0.8em}
  {}

% Subsubsection
\titleformat{\subsubsection}
  {\bfseries\color{navy}}
  {\thesubsubsection}
  {0.6em}
  {}

% ----------------------------------------------------------
%  PAGE STYLE
% ----------------------------------------------------------
\usepackage{fancyhdr}
\pagestyle{fancy}

\fancyhf{} % clear all

\lhead{\textcolor{navy}{\headingfont \courseTitle}}
\rhead{\textcolor{teamred}{\thepage}}
\renewcommand{\headrule}{\color{teamred}\hrule height 1pt}

% ----------------------------------------------------------
%  HYPERLINK STYLE
% ----------------------------------------------------------
\usepackage{hyperref}

\hypersetup{
  colorlinks = true,
  linkcolor  = teamred,
  urlcolor   = navy,
  citecolor  = accentyellow
}

% ----------------------------------------------------------
%  CUSTOM BOX ENVIRONMENTS
% ----------------------------------------------------------
\usepackage{tcolorbox}
\tcbuselibrary{skins,breakable}

% Info box
\newtcolorbox{infobox}{
  enhanced,
  breakable,
  colback   = softwhite,
  colframe  = navy,
  coltext   = navy,
  boxrule   = 1pt,
  arc       = 4pt,
  left=6pt, right=6pt, top=6pt, bottom=6pt
}

% Warning box
\newtcolorbox{warningbox}{
  enhanced,
  breakable,
  colback   = accentyellow!30,
  colframe  = accentyellow!80!navy,
  coltext   = navy,
  boxrule   = 1pt,
  arc       = 4pt,
  left=6pt, right=6pt, top=6pt, bottom=6pt
}

% Highlight box
\newtcolorbox{highlightbox}{
  enhanced,
  breakable,
  colback   = teamred!10,
  colframe  = teamred,
  coltext   = navy,
  boxrule   = 1pt,
  arc       = 4pt,
  left=6pt, right=6pt, top=6pt, bottom=6pt
}

% ----------------------------------------------------------
%  ITEMIZE / ENUMERATE STYLE
% ----------------------------------------------------------
\usepackage{enumitem}

\setlist[itemize]{
  label=\textcolor{teamred}{\large\textbullet},
  leftmargin=1.2em,
  itemsep=4pt
}

\setlist[enumerate]{
  label=\textcolor{navy}{\arabic*.},
  leftmargin=1.4em,
  itemsep=4pt
}

% ----------------------------------------------------------
%  TITLE STYLE
% ----------------------------------------------------------
\usepackage{titling}

\pretitle{
  \begin{center}
  \headingfont\bfseries\LARGE\color{navy}
}
\posttitle{
  \end{center}
}

\preauthor{
  \begin{center}
  \color{teamred}
}
\postauthor{
  \end{center}
}

% ----------------------------------------------------------
%  OPTIONAL — SOFT PAGE BACKGROUND
%  Comment out if you want white pages
% ----------------------------------------------------------
\usepackage{pagecolor}
\pagecolor{cream}
\color{navy}

% ----------------------------------------------------------
% CUSTOM COMMANDS FOR TEMPLATE
% ----------------------------------------------------------
\newcommand{\courseTitle}{Introduction to the Linux Course}
\newcommand{\courseDate}{23/11/2025}

% ----------------------------------------------------------
% END OF THEME
% ----------------------------------------------------------

\begin{document}

% Title with minimal space
\begin{center}
\headingfont\Huge\bfseries\color{navy}
\courseTitle\\[0.5em]
\Large\color{teamred}
\courseDate
\end{center}

\section*{Topics}

\begin{itemize}
    \item \textbf{Linux history and Linux distros}
          \\Link: \href{https://www.youtube.com/watch?v=ShcR4Zfc6Dw}{\texttt{https://www.youtube.com/watch?v=ShcR4Zfc6Dw}}
    
    \item \textbf{Unix philosophy}
          \\Link: \href{https://youtu.be/gojeTqXdBH0?si=kg-iP4KEFjLU-Tto}{\texttt{https://youtu.be/gojeTqXdBH0?si=kg-iP4KEFjLU-Tto}}
    
    \item \textbf{Kernel vs OS}
          \\Link: \href{https://youtu.be/IvGdY6luTtU?si=nEcir9iqjmYMa1lo}{\texttt{https://youtu.be/IvGdY6luTtU?si=nEcir9iqjmYMa1lo}}
    
    \item \textbf{Free Software Philosophy}
          \\Link: \href{https://youtu.be/Ag1AKIl_2GM?si=0YDnS2BDSrAAmFX8}{\texttt{https://youtu.be/Ag1AKIl_2GM?si=0YDnS2BDSrAAmFX8}}
    
    \item \textbf{Why programmers love Linux}
          \\Link: \href{https://youtu.be/otDOHt\_Jges?si=4ChY6-xpYbOdbKQ6}{\texttt{https://youtu.be/otDOHt\_Jges?si=4ChY6-xpYbOdbKQ6}}
    
    \item \textbf{Pros and cons}
          \\Link: \href{https://youtu.be/pYfzZRyRzvs?si=4Q5\_H0gJ3JHBodOc}{\texttt{https://youtu.be/pYfzZRyRzvs?si=4Q5\_H0gJ3JHBodOc}}
\end{itemize}

\vspace{1em}
\begin{center}
\color{teamred}\rule{0.8\textwidth}{1pt}
\end{center}

\section*{Introduction to Linux}

\begin{infobox}
Welcome to the first step of your Linux journey. Whether you're a developer, cybersecurity student, or someone who just wants to understand how computers \textit{really} work, this course will give you the foundation you need to actually \textbf{use} Linux—not just read about it.
\end{infobox}

This introduction will help you understand:

\begin{itemize}
    \item What Linux actually is
    \item Why it became such a big deal
    \item How it's different from other operating systems
    \item Why developers, hackers, and engineers prefer it
    \item What mindset you should adopt while learning Linux
    \item How we'll approach the workshop
\end{itemize}

\begin{warningbox}
If this is your first time touching Linux: \textbf{perfect}. This course is made for you.
\end{warningbox}

\section*{1. Why Linux Matters Today}

If you look at the modern world—servers, cloud platforms, supercomputers, phones, routers, IoT, cybersecurity labs, programming environments—Linux is \textit{everywhere}. It's the invisible engine running most of the technologies you interact with daily.

A few facts:

\begin{itemize}
    \item \textbf{100\%} of the world's top 500 supercomputers run Linux.
    \item \textbf{Most web servers} use Linux (Apache, Nginx).
    \item \textbf{Android}, the most widely-used mobile OS, is built on the Linux kernel.
    \item \textbf{Cybersecurity distributions} (Kali, Parrot) are Linux-based.
    \item Developers choose Linux for transparency, speed, customizability, and control.
\end{itemize}

So even if you've never installed Linux on your laptop, you've definitely used it indirectly.

\section*{2. Understanding Linux: What It Really Is}

One of the biggest misconceptions is thinking Linux is "that thing that replaces Windows." Linux is not a program. It's not a GUI. It's not even an operating system by itself.

\subsection*{Linux = the kernel}

The kernel is the core program that:

\begin{itemize}
    \item Talks to hardware
    \item Manages memory
    \item Schedules processes
    \item Handles permissions
    \item Controls I/O
\end{itemize}

Everything else—the desktop, apps, package manager, and utilities—comes from the broader free-software ecosystem called \textbf{GNU}.

Together:
\textbf{GNU tools + Linux kernel = GNU/Linux (a complete OS).}

This separation is important because Linux teaches you something deeper than using an OS—it teaches you \textit{how systems work}.

\section*{3. Quick Overview of the UNIX Philosophy}

Linux inherits its design principles from Unix. This philosophy will shape how you think during the whole workshop:

\subsection*{The UNIX Philosophy}

\begin{enumerate}
    \item \textbf{Write programs that do one thing, and do it well.}
    \item \textbf{Write programs to work together.}
    \item \textbf{Write programs to handle text, because text is a universal interface.}
\end{enumerate}

This leads to:

\begin{itemize}
    \item Tools that are small and powerful
    \item Commands that do simple tasks
    \item The magic happens when you combine them using pipes (`|`)
\end{itemize}

It's why a Linux power-user is unstoppable—every small tool becomes a Lego block, and you can build anything from those blocks.

\section*{4. Linux History (The Simple Version)}

Let's keep it practical.

\subsection*{1969 – Unix born}
AT&T created Unix. Clean, elegant, powerful.

\subsection*{1983 – GNU Project launched}
Richard Stallman began building a free Unix-like system. GNU had everything—compilers, libraries, tools—\textit{except a kernel.}

\subsection*{1991 – Linus Torvalds writes the Linux Kernel}
As a personal hobby project. He released it under a free license, allowing anyone to study, modify, and contribute.

\subsection*{1992 – GNU + Linux merged}
This created the first fully free operating system.

\subsection*{1990s–2000s – Linux takes over servers}
Why? Stability, control, security, no licensing fees.

\subsection*{2010s–2020s – Linux becomes everywhere}
Cloud, containers, mobile, embedded systems, DevOps… The world quietly standardized on Linux.

\begin{highlightbox}
That's why learning Linux today is a \textbf{career skill}, not just a personal preference.
\end{highlightbox}

\section*{5. Linux Distributions (Distros)}

If Linux is the engine, a \textbf{distro} is the full car.

Different distros = different combinations of:

\begin{itemize}
    \item Package managers
    \item Desktop environments
    \item Default tools
    \item Philosophy
    \item Target audience
\end{itemize}

\subsection*{Beginner-friendly distros}
\begin{itemize}
    \item Ubuntu
    \item Linux Mint
    \item MX Linux
\end{itemize}

\subsection*{Power-user distros}
\begin{itemize}
    \item Arch Linux
    \item Gentoo
    \item Void Linux
\end{itemize}

\subsection*{Enterprise distros}
\begin{itemize}
    \item Red Hat Enterprise Linux
    \item SUSE Linux Enterprise
\end{itemize}

\subsection*{Security distros}
\begin{itemize}
    \item Kali Linux
    \item Parrot OS
\end{itemize}

\begin{infobox}
Throughout the workshop, we'll use a distro that is: Stable, Easy to install, Common across many environments (Usually Ubuntu or Mint—you'll get confirmation before the workshop.)
\end{infobox}

\section*{6. Kernel vs Operating System}

This is where newcomers usually get confused.

\subsection*{Kernel}
\begin{itemize}
    \item The lowest-level component
    \item Manages memory, processes, scheduling
    \item Talks to CPU, RAM, disks
    \item Loads drivers
    \item Handles system calls
\end{itemize}

\subsection*{Operating System}
\begin{itemize}
    \item Kernel
    \item Plus userland tools (GNU utilities)
    \item Plus services
    \item Plus package managers
    \item Plus desktop environments
    \item Plus apps
\end{itemize}

Linux is the \textit{kernel}. GNU/Linux is the \textit{operating system}.

This distinction matters because in this course you will interact with the kernel indirectly using tools like:

\begin{itemize}
    \item \texttt{top}, \texttt{ps}, \texttt{dmesg}
    \item \texttt{/proc} filesystem
    \item System calls through C programs
    \item Permissions and user management
\end{itemize}

This duality is key to becoming a real system administrator or developer.

\section*{7. The Free Software Philosophy (Why Linux Exists)}

Linux exists because of the free software movement. Free software is about \textit{freedom}, not price.

\subsection*{The Four Essential Freedoms}
\begin{enumerate}
    \item \textbf{Freedom 0:} Run the program for any purpose
    \item \textbf{Freedom 1:} Study the source code and modify it
    \item \textbf{Freedom 2:} Distribute copies
    \item \textbf{Freedom 3:} Distribute modified versions
\end{enumerate}

If a program denies any of these freedoms, it becomes \textbf{proprietary}, which means:

\begin{itemize}
    \item You can't see how it works
    \item You can't fix it
    \item You depend on the vendor
    \item You lose control over your own computer
\end{itemize}

Linux was built to solve this.

During the workshop, you'll see the impact of this philosophy:

\begin{itemize}
    \item Software you can fully audit
    \item Tools you can modify
    \item Community-built contributions
    \item Transparency at every layer
\end{itemize}

This is why cybersecurity professionals prefer Linux—you can see \textit{everything} happening in the system.

\section*{8. Why Programmers Love Linux}

Software developers love Linux for reasons that will become obvious once you start using it:

\subsection*{1. Everything is a file}
Devices, sockets, processes, IPC… all represented as files. Simplifies development massively.

\subsection*{2. Powerful terminal environment}
You can automate anything.

\subsection*{3. Package managers}
Install complete development stacks in seconds.

\subsection*{4. DevOps and cloud-native tools}
Docker, Kubernetes, CI pipelines—all built on Linux.

\subsection*{5. Open-source ecosystem}
Learn from real-world code. Fix things. Build tools.

\subsection*{6. High performance and stability}
Perfect for servers and embedded systems.

\subsection*{7. Customizability}
Modify your environment to match exactly how you think and work.

\begin{infobox}
This course will help you feel comfortable enough that Linux becomes your main development environment.
\end{infobox}

\section*{9. Pros and Cons of Linux}

\subsection*{Pros}
\begin{itemize}
    \item Powerful CLI
    \item Free and open source
    \item Fast, lightweight
    \item Secure and stable
    \item Full control over your system
    \item Great for programming
    \item Great for networking
    \item Massive community support
\end{itemize}

\subsection*{Cons}
\begin{itemize}
    \item Some hardware drivers might be missing
    \item Gaming is improving but still imperfect
    \item Adobe apps not officially available
    \item Certain distros can be difficult to use
    \item Requires learning mindset, not point-and-click
\end{itemize}

But as a developer, engineer, or cybersecurity student, the \textit{pros} outweigh the cons by a huge margin.

\section*{10. How to Think Like a Linux User}

To succeed in this workshop:

\subsection*{Adopt these habits:}
\begin{itemize}
    \item \textbf{Be curious} — try commands, explore folders.
    \item \textbf{Read error messages} — Linux errors are usually helpful.
    \item \textbf{Use man pages} — every command has documentation.
    \item \textbf{Experiment safely} — break things in virtual machines.
    \item \textbf{Think in small tools} — pipes, filters, text processing.
    \item \textbf{Love the terminal} — it's your real power.
\end{itemize}

Linux rewards people who explore.

\section*{11. What You'll Learn in This Course}

This introduction sets the ground for the entire course. By the end of the workshop, you will:

\subsection*{System Skills}
\begin{itemize}
    \item Install and manage Linux
    \item Understand filesystems
    \item Work confidently in the terminal
    \item Manage users, processes, permissions
    \item Control systemd services
    \item Understand the kernel's role
\end{itemize}

\subsection*{Networking Skills}
\begin{itemize}
    \item Configure network interfaces
    \item Use diagnostic tools (ping, traceroute, curl)
    \item Manage firewall rules
\end{itemize}

\subsection*{Programming and Automation}
\begin{itemize}
    \item Shell scripting
    \item Using package managers
    \item Working with logs
    \item Troubleshooting
\end{itemize}

\subsection*{Security Awareness}
\begin{itemize}
    \item Permissions
    \item Privilege escalation basics
    \item Monitoring system behavior
\end{itemize}

\begin{highlightbox}
You'll not just use Linux—you'll \textbf{understand} Linux.
\end{highlightbox}

\section*{Useful Resources}

\subsection*{YouTube Channels}
\begin{itemize}
    \item anaHr — \href{https://www.youtube.com/@anaHr}{\texttt{https://www.youtube.com/@anaHr}}
    \item mohammad besar — \href{https://www.youtube.com/@mmbesar}{\texttt{https://www.youtube.com/@mmbesar}}
    \item alwaqqad — \href{https://www.youtube.com/@alwaqad}{\texttt{https://www.youtube.com/@alwaqad}}
    \item sudostart — \href{https://www.youtube.com/@Sudo\_Start}{\texttt{https://www.youtube.com/@Sudo\_Start}}
    \item linuxtopia — \href{https://www.youtube.com/@linuxtopia}{\texttt{https://www.youtube.com/@linuxtopia}}
    \item  abdulmogeeb alhameed — \href{https://www.youtube.com/@abdulmogeeb}{\texttt{https://www.youtube.com/@abdulmogeeb}}
\end{itemize}

\subsection*{Books}
\begin{itemize}
    \item \textbf{The Art Of UNIX Programming} — \textit{Eric Raymond}
\end{itemize}

% Footer
\vspace{\fill}
\begin{center}
\color{muted}
All rights reserved for \texttt{TogetherTeam} [NOT For Commercial Use] !!\\
contact the author via email \href{mailto:khaledmhfz2004@gmail.com}{\texttt{khaledmhfz2004@gmail.com}}
\end{center}
\end{document}